\documentclass[12pt,a4paper]{article}

% ---- Encoding & Font ----
\usepackage{iftex}
\ifPDFTeX
  \usepackage[utf8]{inputenc}
  \usepackage[T1]{fontenc}
\else
  \usepackage{fontspec}  % XeTeX/LuaTeX (tectonic): Unicode nativo (·, —, ¿, ª, §)
\fi
\usepackage[spanish]{babel}
\usepackage{csquotes}

% ---- Page layout ----
\usepackage[top=2.5cm, bottom=2.5cm, left=2.5cm, right=2.5cm]{geometry}
\usepackage{setspace}
\onehalfspacing

% ---- Math ----
\usepackage{amsmath, amssymb, amsthm}
\usepackage{mathtools}
\usepackage{physics}
\usepackage{esint}  % \oiint

% ---- Graphics ----
\usepackage{graphicx}
\usepackage{tikz}
\usepackage{float}
\usepackage{caption}

% ---- Tables ----
\usepackage{array}
\usepackage{booktabs}
\usepackage{tabularx}
\usepackage{colortbl}

% ---- Colours ----
\usepackage{xcolor}
\definecolor{notebg}{HTML}{FFF3CD}
\definecolor{noteborder}{HTML}{FFC107}
\definecolor{boxred}{HTML}{F8D7DA}
\definecolor{boxredborder}{HTML}{F5C6CB}

% ---- Boxes ----
\usepackage{mdframed}
\usepackage{tcolorbox}
\tcbuselibrary{most}

\newtcolorbox{keybox}[1][]{
  colback=blue!5,
  colframe=blue!70!black,
  arc=4pt,
  boxrule=1pt,
  fonttitle=\bfseries,
  title={#1}
}

\newtcolorbox{warnbox}[1][]{
  colback=boxred,
  colframe=boxredborder,
  coltitle=black!75,
  arc=4pt,
  boxrule=1pt,
  fonttitle=\bfseries,
  title={#1}
}

\newtcolorbox{notebox}[1][]{
  colback=notebg,
  colframe=noteborder,
  coltitle=black!75,
  arc=4pt,
  boxrule=1pt,
  fonttitle=\bfseries,
  title={#1}
}

% ---- Hyperlinks & TOC ----
\usepackage[hidelinks]{hyperref}
\usepackage{bookmark}

% ---- Enumerate ----
\usepackage{enumitem}

% ---- Miscellaneous ----
\usepackage{icomma}
\usepackage{siunitx}
\usepackage{textcomp}
\usepackage[bottom]{footmisc}

% ---- Diagramas (circuitos, gráficas y esquemas) ----
\usepackage[europeanresistors]{circuitikz}
\usepackage{pgfplots}
\pgfplotsset{compat=1.18}
\usetikzlibrary{babel}  % babel español activa `>`; esto evita romper las flechas `->`
% courses/ElecMag2024/Clases/assets/tikzstyles.tex
% ---------------------------------------------------------------------------
% Estilos compartidos para los diagramas del curso Electromagnetismo (2024).
% Fuente única del "look" visual (colores, grosores, escalas) para que todas
% las figuras (TikZ / pgfplots / CircuiTikz) sean consistentes entre clases.
%
% Es una COPIA de courses/Fisica3-2015/Clases/assets/tikzstyles.tex: mismo
% dominio, misma paleta. Copia y no symlink para que cada curso sea
% autocontenido y la paleta de Física III pueda evolucionar aparte.
%
% Uso: la clase debe cargar antes `tikz` y `pgfplots` (y `circuitikz` si la
%      clase tiene circuitos), y luego % courses/ElecMag2024/Clases/assets/tikzstyles.tex
% ---------------------------------------------------------------------------
% Estilos compartidos para los diagramas del curso Electromagnetismo (2024).
% Fuente única del "look" visual (colores, grosores, escalas) para que todas
% las figuras (TikZ / pgfplots / CircuiTikz) sean consistentes entre clases.
%
% Es una COPIA de courses/Fisica3-2015/Clases/assets/tikzstyles.tex: mismo
% dominio, misma paleta. Copia y no symlink para que cada curso sea
% autocontenido y la paleta de Física III pueda evolucionar aparte.
%
% Uso: la clase debe cargar antes `tikz` y `pgfplots` (y `circuitikz` si la
%      clase tiene circuitos), y luego % courses/ElecMag2024/Clases/assets/tikzstyles.tex
% ---------------------------------------------------------------------------
% Estilos compartidos para los diagramas del curso Electromagnetismo (2024).
% Fuente única del "look" visual (colores, grosores, escalas) para que todas
% las figuras (TikZ / pgfplots / CircuiTikz) sean consistentes entre clases.
%
% Es una COPIA de courses/Fisica3-2015/Clases/assets/tikzstyles.tex: mismo
% dominio, misma paleta. Copia y no symlink para que cada curso sea
% autocontenido y la paleta de Física III pueda evolucionar aparte.
%
% Uso: la clase debe cargar antes `tikz` y `pgfplots` (y `circuitikz` si la
%      clase tiene circuitos), y luego \input{../assets/tikzstyles.tex}
% (compilar desde el directorio de la clase para que la ruta relativa resuelva).
% ---------------------------------------------------------------------------

% ---- Paleta (alineada con las cajas del preámbulo: keybox/notebox/warnbox) ----
\definecolor{figblue}{HTML}{1F4E79}   % azul principal (líneas, keybox)
\definecolor{figred}{HTML}{C0392B}    % rojo de énfasis (warnbox)
\definecolor{figamber}{HTML}{B8860B}  % ámbar (notebox)
\definecolor{figgray}{HTML}{555555}   % auxiliares, ejes, guías

% ---- CircuiTikz: escala y grosor de los componentes ----
% Guarda \ifdefined: la electrostática (Clases 1-14) no carga `circuitikz`, y
% sin la guarda el \input revienta con `Undefined control sequence \ctikzset`.
% Cuando el curso llegue a circuitos, el bloque se activa solo.
\ifdefined\ctikzset
  \ctikzset{
    bipoles/thickness=1,
    resistors/scale=0.9,
    capacitors/scale=0.9,
  }
\fi

% ---- TikZ: estilos reutilizables ----
% Nota: las puntas se declaran como `-{Latex[...]}` y no como `->`. La forma
% con llaves NO contiene un `>` literal, así que es inmune al `>` activo que
% introduce babel español (ver CLAUDE.md §7.3.3).
\usetikzlibrary{arrows.meta,calc,angles,quotes}

\tikzset{
  figline/.style={figblue, line width=0.9pt},
  figaux/.style={figgray, dashed, line width=0.6pt},
  % vectores
  figvec/.style={figblue, line width=0.9pt, -{Latex[length=2.2mm,width=1.6mm]}},
  figvecres/.style={figred, line width=1.3pt, -{Latex[length=2.6mm,width=1.9mm]}},
  figvecaux/.style={figgray, line width=0.7pt, -{Latex[length=1.8mm,width=1.3mm]}},
  figaxis/.style={figgray, line width=0.6pt, -{Latex[length=2mm,width=1.4mm]}},
  % cargas puntuales
  figcarga/.style={circle, inner sep=0pt, minimum size=5.4pt, draw=none},
  figpos/.style={figcarga, fill=figred},
  figneg/.style={figcarga, fill=figblue},
  figneutra/.style={figcarga, fill=figgray},
  % etiquetas
  figlbl/.style={font=\small},
  figlblaux/.style={font=\footnotesize, figgray},
}

% ---- pgfplots: estilo base de las gráficas ----
\pgfplotsset{
  emfig/.style={
    width=11cm, height=6.5cm,
    axis lines=left,
    xlabel style={font=\small},
    ylabel style={font=\small},
    tick label style={font=\footnotesize},
    every axis plot/.append style={line width=1pt},
    grid=major,
    grid style={figgray!20},
    legend cell align=left,
    legend style={font=\footnotesize, draw=figgray!40},
  },
}

% (compilar desde el directorio de la clase para que la ruta relativa resuelva).
% ---------------------------------------------------------------------------

% ---- Paleta (alineada con las cajas del preámbulo: keybox/notebox/warnbox) ----
\definecolor{figblue}{HTML}{1F4E79}   % azul principal (líneas, keybox)
\definecolor{figred}{HTML}{C0392B}    % rojo de énfasis (warnbox)
\definecolor{figamber}{HTML}{B8860B}  % ámbar (notebox)
\definecolor{figgray}{HTML}{555555}   % auxiliares, ejes, guías

% ---- CircuiTikz: escala y grosor de los componentes ----
% Guarda \ifdefined: la electrostática (Clases 1-14) no carga `circuitikz`, y
% sin la guarda el \input revienta con `Undefined control sequence \ctikzset`.
% Cuando el curso llegue a circuitos, el bloque se activa solo.
\ifdefined\ctikzset
  \ctikzset{
    bipoles/thickness=1,
    resistors/scale=0.9,
    capacitors/scale=0.9,
  }
\fi

% ---- TikZ: estilos reutilizables ----
% Nota: las puntas se declaran como `-{Latex[...]}` y no como `->`. La forma
% con llaves NO contiene un `>` literal, así que es inmune al `>` activo que
% introduce babel español (ver CLAUDE.md §7.3.3).
\usetikzlibrary{arrows.meta,calc,angles,quotes}

\tikzset{
  figline/.style={figblue, line width=0.9pt},
  figaux/.style={figgray, dashed, line width=0.6pt},
  % vectores
  figvec/.style={figblue, line width=0.9pt, -{Latex[length=2.2mm,width=1.6mm]}},
  figvecres/.style={figred, line width=1.3pt, -{Latex[length=2.6mm,width=1.9mm]}},
  figvecaux/.style={figgray, line width=0.7pt, -{Latex[length=1.8mm,width=1.3mm]}},
  figaxis/.style={figgray, line width=0.6pt, -{Latex[length=2mm,width=1.4mm]}},
  % cargas puntuales
  figcarga/.style={circle, inner sep=0pt, minimum size=5.4pt, draw=none},
  figpos/.style={figcarga, fill=figred},
  figneg/.style={figcarga, fill=figblue},
  figneutra/.style={figcarga, fill=figgray},
  % etiquetas
  figlbl/.style={font=\small},
  figlblaux/.style={font=\footnotesize, figgray},
}

% ---- pgfplots: estilo base de las gráficas ----
\pgfplotsset{
  emfig/.style={
    width=11cm, height=6.5cm,
    axis lines=left,
    xlabel style={font=\small},
    ylabel style={font=\small},
    tick label style={font=\footnotesize},
    every axis plot/.append style={line width=1pt},
    grid=major,
    grid style={figgray!20},
    legend cell align=left,
    legend style={font=\footnotesize, draw=figgray!40},
  },
}

% (compilar desde el directorio de la clase para que la ruta relativa resuelva).
% ---------------------------------------------------------------------------

% ---- Paleta (alineada con las cajas del preámbulo: keybox/notebox/warnbox) ----
\definecolor{figblue}{HTML}{1F4E79}   % azul principal (líneas, keybox)
\definecolor{figred}{HTML}{C0392B}    % rojo de énfasis (warnbox)
\definecolor{figamber}{HTML}{B8860B}  % ámbar (notebox)
\definecolor{figgray}{HTML}{555555}   % auxiliares, ejes, guías

% ---- CircuiTikz: escala y grosor de los componentes ----
% Guarda \ifdefined: la electrostática (Clases 1-14) no carga `circuitikz`, y
% sin la guarda el \input revienta con `Undefined control sequence \ctikzset`.
% Cuando el curso llegue a circuitos, el bloque se activa solo.
\ifdefined\ctikzset
  \ctikzset{
    bipoles/thickness=1,
    resistors/scale=0.9,
    capacitors/scale=0.9,
  }
\fi

% ---- TikZ: estilos reutilizables ----
% Nota: las puntas se declaran como `-{Latex[...]}` y no como `->`. La forma
% con llaves NO contiene un `>` literal, así que es inmune al `>` activo que
% introduce babel español (ver CLAUDE.md §7.3.3).
\usetikzlibrary{arrows.meta,calc,angles,quotes}

\tikzset{
  figline/.style={figblue, line width=0.9pt},
  figaux/.style={figgray, dashed, line width=0.6pt},
  % vectores
  figvec/.style={figblue, line width=0.9pt, -{Latex[length=2.2mm,width=1.6mm]}},
  figvecres/.style={figred, line width=1.3pt, -{Latex[length=2.6mm,width=1.9mm]}},
  figvecaux/.style={figgray, line width=0.7pt, -{Latex[length=1.8mm,width=1.3mm]}},
  figaxis/.style={figgray, line width=0.6pt, -{Latex[length=2mm,width=1.4mm]}},
  % cargas puntuales
  figcarga/.style={circle, inner sep=0pt, minimum size=5.4pt, draw=none},
  figpos/.style={figcarga, fill=figred},
  figneg/.style={figcarga, fill=figblue},
  figneutra/.style={figcarga, fill=figgray},
  % etiquetas
  figlbl/.style={font=\small},
  figlblaux/.style={font=\footnotesize, figgray},
}

% ---- pgfplots: estilo base de las gráficas ----
\pgfplotsset{
  emfig/.style={
    width=11cm, height=6.5cm,
    axis lines=left,
    xlabel style={font=\small},
    ylabel style={font=\small},
    tick label style={font=\footnotesize},
    every axis plot/.append style={line width=1pt},
    grid=major,
    grid style={figgray!20},
    legend cell align=left,
    legend style={font=\footnotesize, draw=figgray!40},
  },
}


% ---- Title ----
\title{\textbf{Clase 25: Energía magnética} \\[0.3em]
        \large{Energía de un conjunto de corrientes, densidad de energía y fuerzas sobre circuitos}}
\author{Transcripto y expandido --- Curso Electromagnetismo (OpenFING)}
\date{}

\begin{document}

\maketitle
\thispagestyle{empty}
\newpage

\tableofcontents
\newpage

\section{Almacenar energía en un campo que no trabaja}

La idea parece rara de entrada. La potencia que una fuerza entrega a una carga
es $\vec F\cdot\vec v$, y la fuerza magnética es ortogonal a la velocidad: \textbf{la
fuerza magnética no trabaja}. ¿Cómo se almacena energía en un campo magnético?

Por la \textbf{inducción}: un campo magnético variable genera un campo eléctrico
inducido, y ése sí trabaja y puede dar o tomar energía de las cargas. Por eso,
como se sabe de Física III, para almacenar energía en una bobina hace falta que
las corrientes \textbf{cambien en el tiempo} durante el proceso.

Así como el condensador es el dispositivo típico para almacenar energía en un
campo eléctrico, la \textbf{bobina} o \textbf{inductor} lo es para el magnético. Pero
cualquier distribución de corrientes genera un campo y tiene una energía
asociada; la bobina es «particularmente llamadora» porque con muchas vueltas
concentra mucho campo en una región. Una sola espira almacena menos, con un
campo más desparramado.

\section{Energía en una bobina}

\subsection{Balance con una batería}

\begin{figure}[H]
  \centering
  % Clase 25 — una batería alimenta una bobina real, representada como una
% bobina ideal L en serie con su resistencia interna R. El trabajo de la
% batería se reparte entre el efecto Joule (R I^2 dt) y la energía acumulada
% en la bobina, d(L I^2 / 2).
\begin{tikzpicture}[scale=1]
  \begin{scope}[line width=0.9pt, color=figblue]
    \draw (0,0) to[battery1] (0,2.2);
    \draw (0,2.2) -- (0.8,2.2) to[R] (2.8,2.2) -- (3.6,2.2);
    \draw (3.6,2.2) to[L] (3.6,0);
    \draw (3.6,0) -- (0,0);
  \end{scope}
  \node[figlbl, anchor=east] at (-0.45,1.1) {$V$};
  \node[figlbl, figblue] at (1.8,2.65) {$R$};
  \node[figlbl, figblue, anchor=west] at (3.95,1.1) {$L$};
  \draw[figvecres] (2.2,-0.4) -- (1.2,-0.4);
  \node[figlbl, text=figred, anchor=north] at (1.7,-0.45) {$I$};
\end{tikzpicture}

  \caption{Una batería alimenta una bobina real, representada como una bobina ideal $L$ en serie con su resistencia interna $R$.}
\end{figure}

Una batería alimenta una bobina, representada como una bobina ideal de
autoinductancia $L$ en serie con su resistencia interna $R$. El trabajo de la
batería al mover una carga $dq$ es $V\,dq$, y la diferencia de potencial es la
caída en la resistencia más la FEM de la bobina, que por Faraday es
$d\Phi/dt$. Si la bobina es \textbf{rígida y lineal}, $d\Phi/dt = L\,dI/dt$:

\[
V = R\,I + L\,\frac{dI}{dt}.
\]

La energía por unidad de carga no depende sólo de la corriente, sino de su
derivada. Con $dq = I\,dt$:

\[
dW_b = R I^2\,dt + L I\,\frac{dI}{dt}\,dt
= R I^2\,dt + d\!\left(\tfrac12 L I^2\right).
\]

\begin{itemize}
  \item $R I^2\,dt$: lo que se disipa por \textbf{efecto Joule}, ya conocido.
  \item $d(\tfrac12 L I^2)$: lo que se \textbf{acumula en la bobina}.
\end{itemize}

\begin{keybox}
\[
\boxed{\ U = \tfrac12 L I^2\ }
\]
\end{keybox}

\subsection{Forma general y forma lineal}

El término de acumulación también se escribe como $I\,d\Phi$:

\[
dW_b^{\text{(acum)}} = I\,d\Phi .
\]

\begin{center}
\renewcommand{\arraystretch}{1.35}
\small
\begin{tabularx}{\textwidth}{l>{\raggedright\arraybackslash}X}
\toprule
\textbf{Expresión} & \textbf{Validez} \\
\midrule
$I\,d\Phi$ & \textbf{General}: sale de Faraday sin suponer autoinductancia ni linealidad; vale, por ejemplo, con un núcleo no lineal \\
$d(\tfrac12 L I^2)$ & Sólo en el caso \textbf{lineal}: al derivar se trató a $L$ como constante. Si $L$ dependiera de $I$ habría que derivarla también \\
\bottomrule
\end{tabularx}
\end{center}

La energía aportada a la bobina puede ir a dos lugares: a \textbf{energía almacenada}
o, en medios no lineales, a \textbf{pérdidas por histéresis}, un proceso
irreversible que depende de la historia. Eso «desborda los propósitos del
curso»: de acá en más se supone \textbf{medio lineal} (vacío, materiales no
magnéticos o hierro dulce, aproximadamente lineal).

\section{Energía de un conjunto de corrientes}

\subsection{El trabajo de las baterías}

Varias bobinas con corrientes $I_1, I_2, I_3,\dots$ Se ignora el efecto Joule
---siempre está, pero sólo se arrastra---: baterías ideales. El trabajo es la suma
sobre todas las bobinas, y el flujo en cada una se descompone en los flujos
generados por todas:

\[
dW_b = \sum_i I_i\,d\Phi_i = \sum_{i,j} I_i\,d\Phi_{ij}.
\]

Hasta acá es general. Para \textbf{bobinas rígidas} lo único que cambia es la
corriente, y para \textbf{medios lineales} $d\Phi_{ij} = M_{ij}\,dI_j$ con $M_{ij}$
constante:

\[
dW_b = \sum_{i,j} M_{ij}\,I_i\,dI_j .
\]

\begin{notebox}[Por qué importa la linealidad]
Si el medio es lineal, la energía está
bien definida y no depende del camino por el que se llegó a las corrientes
finales: se puede elegir el que se quiera. Con histéresis, los $M$ dependen
de las corrientes, el trabajo \textbf{depende del camino} y hay ciclos
irreversibles, «como si hubiera un rozamiento» cuando las zonas del material
se alinean con el campo.
\end{notebox}

\subsection{Un camino cómodo: subir todas las corrientes a la vez}

Como en el caso de los conductores cargados, se hacen crecer \textbf{todas las
corrientes proporcionalmente}:

\[
I'_i = \alpha\,I_i,\qquad \alpha: 0\to1,\qquad dI'_j = I_j\,d\alpha .
\]

\[
U = \int_0^1 \sum_{i,j} M_{ij}\,(\alpha I_i)\,(I_j\,d\alpha)
= \sum_{i,j} M_{ij}\,I_i I_j\int_0^1\alpha\,d\alpha
\]

\begin{keybox}
\[
\boxed{\ U = \frac12\sum_{i,j} M_{ij}\,I_i\,I_j\ }
\]
\end{keybox}

Es la energía almacenada en un conjunto de \textbf{bobinas rígidas en medios
lineales}, interactuando entre sí. Si no fueran rígidas, parte del trabajo iría
a deformarlas (trabajo mecánico); si no fueran lineales, parte iría al proceso
irreversible de magnetización. Una magnetización lineal sí está incluida.

\subsection{Ejemplos: una y dos bobinas}

\textbf{Una bobina.} Un solo término, $M_{11} = L$: se recupera $U = \tfrac12 L I^2$.

\textbf{Dos bobinas.} Con $M_{11} = L_1$, $M_{22} = L_2$ y los dos términos cruzados,
que son iguales porque $M_{12} = M_{21} = M$:

\begin{keybox}
\[
\boxed{\ U = \tfrac12 L_1 I_1^2 + \tfrac12 L_2 I_2^2 + M\,I_1 I_2\ }
\]
\end{keybox}

\subsection{Prueba de la cota de la inductancia mutua}

La cota enunciada la clase pasada sale de un argumento físico: con las
corrientes apagadas no hay energía almacenada (nivel cero), y para llegar a
cualquier par de corrientes hubo que \textbf{aportar} energía. Entonces

\[
U(I_1, I_2)\ \ge\ 0 \qquad\text{para todo } I_1, I_2 .
\]

\begin{notebox}
Esta expresión \textbf{no es manifiestamente positiva}: con $M$ arbitraria podría
ser negativa. Que tenga que serlo es la propiedad física, y de ahí sale la
cota.
\end{notebox}

\begin{figure}[H]
  \centering
  % Clase 25 — la energía de dos bobinas como función de I1 con I2 fijo,
% f(I1) = L1 I1^2/2 + M I1 I2 + L2 I2^2/2: una parábola cuyo mínimo está en
% I1 = -M I2 / L1 y vale (I2^2/2)(L2 - M^2/L1). Que el mínimo sea >= 0 es la
% cota M^2 <= L1 L2. Valores ilustrativos: L1 = 1, L2 = 1, I2 = 1 y tres M.
\begin{tikzpicture}
  \begin{axis}[
    emfig, width=10cm, height=6.2cm,
    xmin=-2.4, xmax=1.4, ymin=-0.3, ymax=2.1,
    xtick={-2,-1,0,1}, ytick={0,0.5,1,1.5,2},
    xlabel={$I_1/I_2$}, ylabel={$U/(L_2 I_2^2)$},
    legend style={at={(0.5,-0.4)}, anchor=north, legend columns=3,
      /tikz/every even column/.append style={column sep=1em}},
  ]
    \addplot[figaux, forget plot, domain=-2.4:1.4, samples=2] {0};
    \addplot[figblue!45, domain=-2.4:1.4, samples=80] {0.5*x^2 + 0.3*x + 0.5};
    \addlegendentry{$M^2 = 0{,}09\,L_1L_2$}
    \addplot[figblue!75, domain=-2.4:1.4, samples=80] {0.5*x^2 + 0.7*x + 0.5};
    \addlegendentry{$M^2 = 0{,}49\,L_1L_2$}
    \addplot[figblue, domain=-2.4:1.4, samples=80] {0.5*x^2 + 1.0*x + 0.5};
    \addlegendentry{$M^2 = L_1L_2$}
  \end{axis}
\end{tikzpicture}

  \caption{Energía de dos bobinas en función de $I_1$, con $I_2$ fijo y $L_1 = L_2$, para tres valores de la mutua. El mínimo de cada parábola vale $\tfrac{I_2^2}{2}(L_2 - M^2/L_1)$: que no sea negativo es la cota $M^2\le L_1L_2$, alcanzada con acoplamiento total.}
\end{figure}

Se fija $I_2$ y se mira $f(I_1) = U(I_1, I_2)$. Si $U\ge0$ variando todo, en
particular el mínimo de $f$ es $\ge0$. Derivando:

\[
f'(I_1) = L_1 I_1 + M I_2 = 0
\quad\Longrightarrow\quad
I_1^{\min} = -\frac{M}{L_1}\,I_2 .
\]

Sustituyendo, todos los términos quedan proporcionales a $I_2^2$:

\[
f(I_1^{\min}) = I_2^2\left(\frac{M^2}{2L_1} + \frac{L_2}{2} - \frac{M^2}{L_1}\right)
= \frac{I_2^2}{2}\left(L_2 - \frac{M^2}{L_1}\right)\ \ge\ 0 .
\]

Como $L_1 > 0$:

\begin{keybox}
\[
\boxed{\ M^2 \le L_1\,L_2\ }
\]
\end{keybox}

La igualdad corresponde al acoplamiento total (los dos bobinados superpuestos
de la clase pasada); con bobinas alejadas $M$ puede ser cero, y hay todos los
casos intermedios.

\section{Densidad de energía magnética}

Sección 9.2. Igual que en el caso eléctrico la energía de los conductores
cargados podía pensarse almacenada en las cargas \textbf{o en el campo}, acá puede
pensarse almacenada en las corrientes o \textbf{en el campo magnético}. Se reescribe
$U$ como integral de volumen de una densidad.

Hipótesis: \textbf{medio lineal} y \textbf{corrientes localizadas} ---hay una esfera finita
fuera de la cual no hay corrientes---, como en el desarrollo multipolar. Las
situaciones realistas son así: «no tenemos ningún cable que nos permita llevar
energía al fin del universo».

\subsection{De los flujos al potencial vector}

Como $\Phi_i = \sum_j M_{ij} I_j$,

\[
U = \frac12\sum_i I_i\,\Phi_i .
\]

Y el flujo por la superficie $S_i$ que se apoya en el circuito $C_i$, usando
$\vec B = \curl\vec A$ y el teorema de Stokes, es la circulación del potencial
vector:

\[
\Phi_i = \int_{S_i}(\curl\vec A)\cdot\hat n\,dS = \oint_{C_i}\vec A\cdot d\vec l,
\qquad
U = \frac12\sum_i I_i\oint_{C_i}\vec A\cdot d\vec l .
\]

Con la regla de siempre para pasar de circuitos a corrientes distribuidas,
$I\,d\vec l\to\vec J\,dV$ (imaginar el cable grueso como muchos cablecitos):

\[
U = \frac12\int_V \vec A\cdot\vec J\;dV,
\]
con $V$ un volumen que incluye todas las corrientes.

\begin{notebox}
A la pregunta de si la regla es una igualdad: es una \textbf{equivalencia}, una
regla mnemotécnica; ninguna de las dos expresiones existe sola, siempre van
multiplicadas por algo (acá por $\vec A$).
\end{notebox}

\subsection{De las corrientes a los campos}

Por Ampère en el material, $\vec J = \curl\vec H$, donde $\vec J$ es la
corriente \textbf{de conducción} (sin la de magnetización). Se usa la identidad

\[
\div(\vec A\times\vec H) = \vec H\cdot(\curl\vec A) - \vec A\cdot(\curl\vec H)
\]

\begin{notebox}
El docente la escribe primero con un signo mal y se lo corrigen. Bromea con
que los profesores del práctico están enojados por el uso de Levi-Civita.
\end{notebox}

Despejando $\vec A\cdot\curl\vec H$ y usando $\curl\vec A = \vec B$:

\[
U = \frac12\int_V \vec H\cdot\vec B\;dV
\;-\; \frac12\oiint_S (\vec A\times\vec H)\cdot\hat n\,dS,
\]
con $S$ el borde de $V$. Integrar en una región más grande que la de las
corrientes no cambia nada («quien puede más puede menos»).

El primer término recuerda al eléctrico, $\tfrac12\vec D\cdot\vec E$ (en el
vacío, $\tfrac12\varepsilon_0E^2$), y en el vacío da el $B^2/2\mu_0$ de Física
III. Hay además, como allá, un \textbf{término de borde}.

\subsection{El término de borde se anula}

Como en el caso eléctrico, se toma $S$ una esfera de radio $R\to\infty$. Para
corrientes localizadas, por el desarrollo multipolar:

\begin{center}
\renewcommand{\arraystretch}{1.35}
\small
\begin{tabularx}{\textwidth}{l>{\raggedright\arraybackslash}X>{\raggedright\arraybackslash}X}
\toprule
\textbf{Cantidad} & \textbf{Decaimiento} & \textbf{Por qué} \\
\midrule
$|\vec A|$ & $1/R^2$ & No hay término monopolar \\
$|\vec B|$, $|\vec H|$ & $1/R^3$ & Rotacional de $\vec A$; $\vec H\propto\vec B$ en medio lineal \\
Área de $S$ & $R^2$ &  \\
Término de borde & $1/R^3$ & $\to 0$ \\
\bottomrule
\end{tabularx}
\end{center}

Tiende a cero, y más rápido que en el caso eléctrico, donde el potencial iba
como $1/R$ y $\vec D$ como $1/R^2$, y el término de borde decrecía como $1/R$.
Integrando sobre todo el espacio:

\begin{keybox}
\[
\boxed{\ U = \frac12\int_{\text{todo el espacio}}\vec H\cdot\vec B\;dV,
\qquad u_B = \frac12\,\vec H\cdot\vec B\ }
\]
\end{keybox}

Vale también con varias zonas lineales distintas (vacío, $\mu_1$, $\mu_2$\dots{}),
usando en cada una su permeabilidad. Con un único medio, $\vec B = \mu\vec H$:

\[
u_B = \frac{B^2}{2\mu},
\]
que en el vacío es la fórmula de Física III.

\begin{notebox}[Regla mnemotécnica]
Del caso eléctrico al magnético, $\vec E\to\vec B$ y
$\vec D\to\vec H$: de $\tfrac12\vec E\cdot\vec D$ a $\tfrac12\vec B\cdot\vec H$.
\end{notebox}

\begin{notebox}
En medios no lineales la noción misma de energía magnética es complicada:
hay una energía potencial asociada, pero también disipación por histéresis,
y el trabajo depende del camino.
\end{notebox}

\section{Fuerzas y torques sobre circuitos rígidos}

Sección 9.3. Como en el caso eléctrico (un dieléctrico que entra en un
condensador), se calculan fuerzas y torques a partir de la energía cuando una
parte del sistema se puede mover respecto de otra. Ejemplo típico: un solenoide
con un núcleo de \textbf{hierro dulce} (para que sea aproximadamente lineal)
parcialmente adentro. El sistema ya no es rígido, pero lo visto se aplica.

\subsection{Balance de energía}

El trabajo mecánico en un desplazamiento $d\vec l$ de una parte respecto de la
otra sale de dos aportes: la batería y la energía almacenada.

\[
\vec F\cdot d\vec l = dW_b - dU .
\]

\subsection{Sistema aislado: un caso poco interesante}

Sin batería, $dW_b = 0$ y $\vec F = -\grad U$, como en el caso eléctrico
aislado. Pero acá \textbf{no es interesante}: sin batería, por el efecto Joule, las
corrientes tienden a cero muy rápido. En el caso eléctrico las cargas
electrostáticas no se disipan, y las fuerzas sobreviven sin que nadie aporte
energía.

\subsection{Corrientes constantes}

El caso realista: las baterías aportan energía de modo que \textbf{las corrientes se
mantienen constantes} (el análogo de las baterías que mantenían los potenciales
fijos en el caso eléctrico). Entonces sólo varían los flujos:

\[
dW_b = \sum_i I_i\,d\Phi_i,
\qquad
dU = \frac12\sum_i I_i\,d\Phi_i
\quad\Longrightarrow\quad
\boxed{\ dW_b = 2\,dU\ }
\]

Igual que en el caso eléctrico a potencial fijo, la batería aporta el doble de
lo que varía la energía almacenada. Entonces $\vec F\cdot d\vec l = 2dU - dU = dU$:

\begin{keybox}
\[
\boxed{\ \vec F = +\left(\grad U\right)_{I},
\qquad \tau_n = +\left(\pdv{U}{\theta}\right)_{I}\ }
\]
\end{keybox}
con $\tau_n$ el torque según un eje $\hat n$ y $\theta$ el ángulo de rotación
alrededor de él. Siempre en medio lineal.

\begin{notebox}[El signo «raro»]
En mecánica la fuerza es menos el gradiente de la
energía potencial cuando ésa es la única fuente de energía. Acá hay dos
aportes, y la batería, al mantener fijas las corrientes, compensa justo el
signo.
\end{notebox}

\subsection{Ejemplo: núcleo de hierro dulce que entra en un solenoide}

Un solenoide ideal de largo $\ell$, sección $A$ y $N$ vueltas, con corriente
$I$. Un núcleo de hierro dulce lineal, de permeabilidad $\mu$, ocupa un tramo de
largo $x$; el resto, de largo $\ell - x$, es vacío.

\begin{notebox}
El docente cambia a mitad de camino las $L$ mayúsculas de los largos por
minúsculas, para no confundirlas con autoinductancias.
\end{notebox}

\begin{figure}[H]
  \centering
  % Clase 25 — solenoide de largo l, sección A y N vueltas con corriente I, con
% un núcleo de hierro dulce (mu) que ocupa un largo x; el resto es vacío.
% A corriente constante, F_x = dU/dx = (I^2 N^2 A / 2 l^2)(mu - mu0) > 0:
% el núcleo es atraído hacia adentro.
\begin{tikzpicture}[scale=1]

  \fill[figgray!8] (0,-0.8) rectangle (7,0.8);
  % nucleo
  \fill[figamber!35] (3.8,-0.62) rectangle (8.6,0.62);
  \draw[figamber, line width=0.9pt] (3.8,-0.62) rectangle (8.6,0.62);
  \node[figlbl] at (8.0,0) {$\mu$};
  \node[figlbl] at (1.9,0) {$\mu_0$};
  % solenoide
  \draw[figline] (0,0.8) -- (7,0.8) (0,-0.8) -- (7,-0.8);
  \foreach \xx in {0.2,0.6,...,6.8} {
    \draw[figblue, line width=1pt] (\xx,-0.86) -- (\xx+0.14,0.86);
  }
  \node[figlbl, text=figblue, anchor=east] at (-0.15,0.55) {$N,\ I$};

  % cotas
  \draw[figgray, line width=0.6pt, {Latex[length=1.6mm]}-{Latex[length=1.6mm]}] (0,-1.2) -- (7,-1.2);
  \node[figlbl, anchor=north] at (3.5,-1.25) {$\ell$};
  \draw[figgray, line width=0.6pt, {Latex[length=1.6mm]}-{Latex[length=1.6mm]}] (3.8,1.15) -- (7,1.15);
  \node[figlbl, anchor=south] at (5.4,1.2) {$x$};

  % fuerza: hacia x creciente, es decir, hacia adentro del solenoide
  \draw[figvecres] (10.0,0) -- (8.75,0);
  \node[figlbl, text=figred, anchor=south] at (9.4,0.1) {$F_x$};

\end{tikzpicture}

  \caption{Núcleo de hierro dulce que ocupa un largo $x$ dentro de un solenoide. A corriente constante, la fuerza va hacia $x$ creciente: el núcleo es atraído hacia adentro.}
\end{figure}

\textbf{Aproximación.} Si $x$ y $\ell - x$ son mucho mayores que el radio del
solenoide, se desprecian los efectos de borde ---dos veces: al usar el campo del
solenoide ideal y al calcular la energía---. Las zonas mal descritas son chicas,
aportan poco, y además lo que aportan casi no depende de $x$.

En cada zona $H = N I/\ell$ (las vueltas por unidad de largo son las mismas), y

\[
B_0 = \mu_0\frac{N}{\ell}I,\qquad B_1 = \mu\frac{N}{\ell}I .
\]

Con $u = B^2/2\mu$ por el volumen de cada zona:

\[
U = \frac{B_0^2}{2\mu_0}\,A\,(\ell-x) + \frac{B_1^2}{2\mu}\,A\,x
= \frac{I^2 N^2 A}{2\ell^2}\Big[\mu_0\,(\ell-x) + \mu\,x\Big].
\]

A corriente constante, la fuerza en la dirección de $x$ creciente:

\begin{keybox}
\[
\boxed{\ F_x = \pdv{U}{x} = \frac{I^2 N^2 A}{2\ell^2}\,(\mu-\mu_0)\ }
\]
\end{keybox}

Como $\mu > \mu_0$, $F_x > 0$: \textbf{el núcleo es atraído hacia adentro} del
solenoide.

\noindent\emph{La próxima clase: circuitos con resistencias, inductancias y condensadores.}

\end{document}
